\documentclass[11pt]{article}

\usepackage{amsmath, amssymb, amsthm}
\usepackage{mathtools}
\usepackage[colorlinks=true]{hyperref}
\usepackage[capitalize]{cleveref}
\usepackage{tikz}
\usetikzlibrary{cd}
\usepackage{pdflscape}

\theoremstyle{definition}
\newtheorem{definition}{Definition}[section]
\newtheorem{theorem}[definition]{Theorem}
\newtheorem{proposition}[definition]{Proposition}
\newtheorem{lemma}[definition]{Lemma}
\newtheorem{remark}[definition]{Remark}

\newcommand{\id}{\mathrm{id}}
\newcommand{\outp}[1]{#1^{+}}
\newcommand{\inpt}[1]{#1^{-}}
\newcommand{\lensob}[1]{\binom{\inpt{#1}}{\outp{#1}}}
\newcommand{\cat}[1]{\mathcal{#1}}
\newcommand{\Cat}[1]{\mathbf{#1}}
\newcommand{\Fun}[1]{\textsf{#1}}
\newcommand{\Lens}[1]{\Cat{Lens}_{#1}}
\newcommand{\LensP}[2]{\Cat{Lens}_{#1}^{#2}}
\newcommand{\op}{^{\mathrm{op}}}
\newcommand{\ev}{\mathrm{ev}}
\newcommand{\coev}{\mathrm{coev}}
\newcommand{\To}[1]{\xrightarrow{#1}}
\newcommand{\ihom}[2]{[#1,#2]}
\newcommand{\smset}{\Cat{Set}}
\newcommand{\qqand}{\qquad\text{and}\qquad}
\DeclareMathSymbol\fatsemi\mathop{AMSa}{"3B}
\newcommand{\then}{\mathbin{\fatsemi}}

\begin{document}

\section{Monoidal promonads and their lens categories}\label{sec.promonad_theta}

Throughout, let $(\cat{C},I,\otimes)$ be a symmetric monoidal closed category with internal hom $\ihom{-}{-}$, evaluation $\ev\colon\ihom{A}{B}\otimes A\to B$, and coevaluation $\coev\colon X\to\ihom{A}{A\otimes X}$.

\subsection{Promonads}

\begin{definition}[Promonad]\label{def.promonad}
A \emph{promonad} on $\cat{C}$ is a profunctor $P\colon\cat{C}\op\times\cat{C}\to\smset$ equipped with
\begin{itemize}
\item a \emph{unit} $\eta_{A,B}\colon\cat{C}(A,B)\to P(A,B)$, natural in $A,B$, and
\item a \emph{multiplication} $\mu_{A,C}\colon\int^{B}P(A,B)\times P(B,C)\to P(A,C)$, natural in $A,C$,
\end{itemize}
satisfying associativity and unit axioms. We write $q\bullet p\coloneqq\mu(p,q)$ for the composite of $p\in P(A,B)$ and $q\in P(B,C)$, so that $q\bullet p\in P(A,C)$.

Equivalently, a promonad is a category $\cat{K}$ with $\mathrm{Ob}(\cat{K})=\mathrm{Ob}(\cat{C})$ and an identity-on-objects functor $J\colon\cat{C}\to\cat{K}$, where $P(A,B)=\cat{K}(A,B)$.
\end{definition}

\begin{remark}[Functorial actions]\label{rmk.functorial}
Since $P$ is a profunctor, it has contravariant and covariant functorial actions. For $h\colon A'\to A$ in $\cat{C}$ and $p\in P(A,B)$, write $P(h,\id)(p)=p\bullet\eta(h)\in P(A',B)$ (precomposition). For $f\colon B\to B'$ and $p\in P(A,B)$, write $P(\id,f)(p)=\eta(f)\bullet p\in P(A,B')$ (postcomposition). The second equality in each case is by the unit axiom of the promonad.
\end{remark}

\subsection{Lax monoidal promonads}

\begin{definition}[Lax monoidal promonad]\label{def.monoidal_promonad}
A promonad $P$ on $(\cat{C},I,\otimes)$ is \emph{lax monoidal} if $P\colon\cat{C}\op\times\cat{C}\to\smset$ is lax monoidal with respect to the external monoidal structure $(A,B)\otimes(A',B')=(A\otimes A',B\otimes B')$ on $\cat{C}\op\times\cat{C}$ and $(\smset,1,\times)$. That is, there are natural maps
\[
\lambda\colon P(A,B)\times P(A',B')\to P(A\otimes A',B\otimes B')
\qqand
\lambda_0\colon\ast\to P(I,I)
\]
satisfying associativity, unit, and the following compatibility conditions:
\begin{enumerate}
\item \textbf{Lax-unit compatibility:} $\lambda(\eta(f),\eta(g))=\eta(f\otimes g)$ for $f\in\cat{C}(A,B)$, $g\in\cat{C}(A',B')$.
\item \textbf{Interchange law:} $\lambda(q\bullet p,\,s\bullet r)=\lambda(q,s)\bullet\lambda(p,r)$ for $p\in P(A,B)$, $q\in P(B,C)$, $r\in P(A',B')$, $s\in P(B',C')$.
\end{enumerate}
Equivalently, $\cat{K}$ is a symmetric monoidal category with $J\colon\cat{C}\to\cat{K}$ strong monoidal. We write $p\mathbin{\hat\otimes} p'\coloneqq\lambda(p,p')$.
\end{definition}

\begin{remark}[Strength]\label{rmk.tambara}
A lax monoidal promonad has a derived \emph{strength}: for $p\in P(A,B)$ and any object $M$,
\[
\lambda(\eta(\id_M),p)\in P(M\otimes A,M\otimes B).
\]
By lax-unit compatibility, $\lambda(\eta(\id_M),\eta(f))=\eta(\id_M\otimes f)$. By the interchange law, $\lambda(\eta(\id_M),-)$ preserves composition:
\begin{equation}\label{eqn.tambara_compose}
\lambda(\eta(\id_M),q\bullet p)=\lambda(\eta(\id_M),q)\bullet\lambda(\eta(\id_M),p).
\end{equation}
\end{remark}

\begin{lemma}[Naturality of $\lambda(\eta(\id),-)$]\label{lem.tambara_nat}
For $p\in P(A,B)$ and $\psi\colon M\to M'$ in $\cat{C}$:
\begin{equation}\label{eqn.tambara_nat}
P(\id,\psi\otimes\id_B)\bigl(\lambda(\eta(\id_M),p)\bigr)=P(\psi\otimes\id_A,\id)\bigl(\lambda(\eta(\id_{M'}),p)\bigr).
\end{equation}
Both sides lie in $P(M\otimes A,M'\otimes B)$.
\end{lemma}
\begin{proof}
Both sides equal $\lambda(\eta(\psi),p)$: the left by naturality of $\lambda$ in the covariant direction, the right by naturality in the contravariant direction.
\end{proof}

\subsection{$P$-algebras}

\begin{definition}[$P$-algebra]\label{def.P_alg}
Let $P$ be a promonad on $\cat{C}$. A \emph{$P$-algebra} on an object $m\in\cat{C}$ is a natural transformation $\alpha\colon P(-,m)\Rightarrow\cat{C}(-,m)$ satisfying:
\begin{enumerate}
\item \textbf{Unit:} $\alpha_X(\eta(f))=f$ for every $f\colon X\to m$ in $\cat{C}$.
\item \textbf{Associativity:} For $p\in P(X,Y)$ and $q\in P(Y,m)$,
\begin{equation}\label{eqn.P_alg_assoc}
\alpha_X(q\bullet p)=\alpha_X\bigl(\eta(\alpha_Y(q))\bullet p\bigr).
\end{equation}
\end{enumerate}
Naturality of $\alpha$ in the first variable means: for $h\colon X'\to X$ in $\cat{C}$,
\begin{equation}\label{eqn.alpha_nat}
\alpha_{X'}(P(h,\id)(q))=\alpha_X(q)\circ h.
\end{equation}
By \cref{rmk.functorial}, the left side is $\alpha_{X'}(q\bullet\eta(h))$, and associativity with the unit axiom recovers the right side.
\end{definition}

\begin{remark}\label{rmk.assoc_push}
By the unit axiom, the right side of \eqref{eqn.P_alg_assoc} simplifies: $\eta(\alpha_Y(q))\bullet p=P(\id,\alpha_Y(q))(p)$, the postcomposition/pushforward of $p$ along the $\cat{C}$-morphism $\alpha_Y(q)$. So associativity says:
\[
\alpha_X(q\bullet p)=\alpha_X\bigl(P(\id,\alpha_Y(q))(p)\bigr).
\]
\end{remark}

\begin{definition}[Monoidal $P$-algebra]\label{def.monoidal_P_alg}
Let $P$ be a lax monoidal promonad on $(\cat{C},I,\otimes)$ and $(m,\mathfrak{m},\mathfrak{e})$ a monoid in $\cat{C}$. A $P$-algebra $\alpha$ on $m$ is \emph{monoidal} if for all $p\in P(X,m)$, $p'\in P(X',m)$:
\begin{align}
\alpha_{X\otimes X'}\bigl(P(\id,\mathfrak{m})(p\mathbin{\hat\otimes} p')\bigr)
  &=\mathfrak{m}\circ(\alpha_X(p)\otimes\alpha_{X'}(p')),\label{eqn.mon_alg_mult}\\
\alpha_I\bigl(P(\id,\mathfrak{e})(\lambda_0)\bigr)&=\mathfrak{e}.\label{eqn.mon_alg_unit}
\end{align}
We call $(m,\mathfrak{m},\mathfrak{e},\alpha)$ a \emph{$P$-monoid}.
\end{definition}

\subsection{The lens category $\LensP{\cat{C}}{P}$}

\begin{definition}\label{def.lensP}
Let $P$ be a lax monoidal promonad on $(\cat{C},I,\otimes)$. The category $\LensP{\cat{C}}{P}$ has:
\begin{itemize}
\item \textbf{Objects:} pairs $\lensob{c}=\binom{\inpt{c}}{\outp{c}}$ where $\outp{c}$ is a comonoid in $\cat{C}$ with comultiplication $\delta_{\outp{c}}$ and counit $\varepsilon_{\outp{c}}$.
\item \textbf{Morphisms $\lensob{c}\to\lensob{d}$:} pairs $(\outp{f},\inpt{f})$ where $\outp{f}\colon\outp{c}\to\outp{d}$ is a comonoid homomorphism and $\inpt{f}\in P(\outp{c}\otimes\inpt{d},\inpt{c})$.
\item \textbf{Identity:} at $\lensob{c}$, the pair $(\id_{\outp{c}},\eta(\varepsilon_{\outp{c}}\otimes\id_{\inpt{c}}))$.
\item \textbf{Composition:} given $(\outp{f},\inpt{f})\colon\lensob{c}\to\lensob{d}$ and $(\outp{g},\inpt{g})\colon\lensob{d}\to\lensob{e}$, their composite has forward part $\outp{g}\circ\outp{f}$ and backward part
\begin{equation}\label{eqn.lens_compose}
\inpt{f}\bullet\lambda(\eta(\id_{\outp{c}}),\inpt{g})\bullet\eta\bigl((\id\otimes\outp{f}\otimes\id)\circ(\delta_{\outp{c}}\otimes\id)\bigr)\;\in\;P(\outp{c}\otimes\inpt{e},\inpt{c}).
\end{equation}
Using \cref{rmk.functorial}, this simplifies to $\inpt{f}\bullet P\bigl((\id\otimes\outp{f}\otimes\id)\circ(\delta\otimes\id),\;\id\bigr)\bigl(\lambda(\eta(\id_{\outp{c}}),\inpt{g})\bigr)$, i.e.\ the composite in $\cat{K}$:
\[
\outp{c}\otimes\inpt{e}\To{\delta\otimes\id}\outp{c}\otimes\outp{c}\otimes\inpt{e}\To{\id\otimes\outp{f}\otimes\id}\outp{c}\otimes\outp{d}\otimes\inpt{e}\To{\lambda(\eta(\id_{\outp{c}}),\inpt{g})}\outp{c}\otimes\inpt{d}\To{\inpt{f}}\inpt{c}
\]
where the first two arrows are $C$-morphisms (via $\eta$) and the last two are $\cat{K}$-morphisms.
\end{itemize}
When $P$ is lax monoidal, $\LensP{\cat{C}}{P}$ inherits a monoidal structure with $\lensob{c}\otimes\lensob{d}=\binom{\inpt{c}\otimes\inpt{d}}{\outp{c}\otimes\outp{d}}$.
\end{definition}

\begin{remark}
When $P=P_{\Fun{T}}$ is the Kleisli profunctor $P_{\Fun{T}}(A,B)=\cat{C}(A,\Fun{T}B)$ of a strong monad $\Fun{T}$, the category $\LensP{\cat{C}}{P_{\Fun{T}}}$ is precisely the coKleisli category $\Lens{\cat{C}}^{\Fun{T}}$ of \cite[Prop.~3.7]{dynamic-algebra-potentials}.
\end{remark}

\subsection{Moore internalization}

\begin{definition}[The map $\phi$]\label{def.phi}
Let $(m,\alpha)$ be a $P$-algebra in monoidal closed $\cat{C}$. Write $H_X\coloneqq\ihom{\inpt{X}}{m}$ for each lens object $\lensob{X}$. Given a morphism $(\outp{f},\inpt{f})\colon\lensob{c}\to\lensob{d}$ in $\LensP{\cat{C}}{P}$, define the $\cat{C}$-morphism
\begin{equation}\label{eqn.phi}
\phi_f\coloneqq\alpha_{\outp{c}\otimes\inpt{d}\otimes H_c}\bigl(P(\id,\ev)\bigl(\lambda(\eta(\id_{H_c}),\inpt{f})\bigr)\bigr)\;\colon\;\outp{c}\otimes\inpt{d}\otimes H_c\to m
\end{equation}
constructed in three steps:
\begin{enumerate}
\item \textbf{Strength:} $\lambda(\eta(\id_{H_c}),\inpt{f})\in P(\outp{c}\otimes\inpt{d}\otimes H_c,\;\inpt{c}\otimes H_c)$,
\item \textbf{Pushforward by $\ev$:} $P(\id,\ev)\bigl(\lambda(\eta(\id_{H_c}),\inpt{f})\bigr)\in P(\outp{c}\otimes\inpt{d}\otimes H_c,\;m)$,
\item \textbf{Drop:} apply $\alpha$ to obtain $\phi_f\colon\outp{c}\otimes\inpt{d}\otimes H_c\to m$ in $\cat{C}$.
\end{enumerate}
Write $\psi_f\colon\outp{c}\otimes H_c\to H_d$ for the adjoint transpose, i.e.\ $\phi_f=\ev\circ(\id_{\inpt{d}}\otimes\psi_f)$ after swapping the second and third tensor factors.
\end{definition}

\begin{definition}[Moore internalization $\Theta_m$]\label{def.Theta}
Define $\Theta_m\colon\LensP{\cat{C}}{P}\to\cat{C}$ on objects by $\lensob{c}\mapsto\outp{c}\otimes H_c$ and on a morphism $(\outp{f},\inpt{f})\colon\lensob{c}\to\lensob{d}$ by the composite
\begin{equation}\label{eqn.Theta_mor}
\outp{c}\otimes H_c\To{\delta\otimes\id}\outp{c}\otimes\outp{c}\otimes H_c\To{\outp{f}\otimes\psi_f}\outp{d}\otimes H_d.
\end{equation}
Equivalently, writing $\coev_{\inpt{d}}'\colon X\to\ihom{\inpt{d}}{\outp{c}\otimes\inpt{d}\otimes X}$ for the coevaluation composed with the symmetry $\inpt{d}\otimes\outp{c}\cong\outp{c}\otimes\inpt{d}$ (so that the argument is $+/-$ clean):
\begin{equation}\label{eqn.Theta_chain}
\outp{c}\otimes H_c\To{\delta\otimes\id}{\outp{c}}^2\otimes H_c\To{\outp{f}\otimes\coev_{\inpt{d}}'}\outp{d}\otimes\ihom{\inpt{d}}{\outp{c}\otimes\inpt{d}\otimes H_c}\To{\ihom{\inpt{d}}{\phi_f}}\outp{d}\otimes H_d.
\end{equation}
\end{definition}

\begin{remark}[Comparison with the monad case]\label{rmk.comparison}
For a strong monad $\Fun{T}$, the map $\phi_f$ decomposes into five steps applied inside $\ihom{\inpt{d}}{-}$:
\begin{equation}\label{eqn.monad_chain}
\outp{c}\otimes\inpt{d}\otimes H_c\To{\inpt{f}\otimes\id}\Fun{T}\inpt{c}\otimes H_c\To{\eta\otimes\id}\Fun{T}\inpt{c}\otimes\Fun{T}H_c\To{\tau}\Fun{T}(\inpt{c}\otimes H_c)\To{\Fun{T}\ev}\Fun{T}m\To{a}m.
\end{equation}
For a general promonad, $\inpt{f}\in P(\outp{c}\otimes\inpt{d},\inpt{c})$ is not a $\cat{C}$-morphism, so it cannot be placed inside $\ihom{\inpt{d}}{-}$. The three-step construction of \cref{def.phi} (Tambara, pushforward, drop via $\alpha$) collapses \eqref{eqn.monad_chain} into a single $\cat{C}$-morphism $\phi_f$.
\end{remark}


\section{The theorem and proof plan}

\begin{theorem}\label{thm.promonad_theta}
Let $P$ be a lax monoidal promonad on a symmetric monoidal closed category $\cat{C}$, and let $(m,\mathfrak{m},\mathfrak{e},\alpha)$ be a $P$-monoid. Then $\Theta_m\colon\LensP{\cat{C}}{P}\to\cat{C}$ is a functor.

If moreover $(m,\alpha)$ is a monoidal $P$-algebra, then $\Theta_m$ is lax monoidal.
\end{theorem}

\begin{lemma}[$\psi$ composition]\label{lem.psi_compose}
Let $f\colon\lensob{c}\to\lensob{d}$ and $g\colon\lensob{d}\to\lensob{e}$ be morphisms in $\LensP{\cat{C}}{P}$, with $\psi_f$ and $\psi_g$ as in \cref{def.phi}. Then
\begin{equation}\label{eqn.psi_compose}
\psi_{g\circ f}=\psi_g\circ(\outp{f}\otimes\psi_f)\circ(\delta_{\outp{c}}\otimes\id_{H_c}).
\end{equation}
\end{lemma}

\begin{proof}
Write $\mathrm{cm}\coloneqq(\id_{\outp{c}}\otimes\outp{f}\otimes\id_{\inpt{e}})\circ(\delta_{\outp{c}}\otimes\id_{\inpt{e}})$. Both $\cat{K}$-composites from $(\outp{c}\otimes H_c)\otimes\outp{d}\otimes\inpt{e}$ to $m$ in the following diagram have the same $\alpha$-image:
\[
\begin{tikzcd}[column sep=3em, row sep=2.5em, every label/.append style={font=\scriptsize}, ampersand replacement=\&]
(\outp{c}\otimes H_c)\otimes\outp{d}\otimes\inpt{e}
  \ar[r, "{\lambda(\eta(\id_{\outp{c}\otimes H_c}),\,\inpt{g})}"]
  \ar[d, "{\psi_f\otimes\id}"']
\& (\outp{c}\otimes H_c)\otimes\inpt{d}
  \ar[r, "{\lambda(\eta(\id_{H_c}),\,\inpt{f})}"]
  \ar[d, "{\psi_f\otimes\id}"]
\& \inpt{c}\otimes H_c
  \ar[r, "{\eta(\ev)}"]
\& m
  \ar[d, equal]
\\
H_d\otimes\outp{d}\otimes\inpt{e}
  \ar[r, "{\lambda(\eta(\id_{H_d}),\,\inpt{g})}"']
\& H_d\otimes\inpt{d}
  \ar[rr, "{\eta(\ev)}"']
\&\& m
\end{tikzcd}
\]
Left tile: $\lambda$-naturality (\cref{lem.tambara_nat}). Right tile: definition of $\phi_f$ ($\alpha$-associativity, \cref{rmk.assoc_push}).

Precomposing with $\eta(\mathrm{cm}\otimes\id_{H_c})$ and applying $\alpha$: the top path gives $\phi_{g\circ f}$ by the interchange law \eqref{eqn.tambara_compose}; the bottom path gives $\phi_g\circ(\psi_f\otimes\id)\circ(\mathrm{cm}\otimes\id_{H_c})$ by $\alpha$-naturality \eqref{eqn.alpha_nat}. Currying with respect to $\inpt{e}$ yields \eqref{eqn.psi_compose}.
\end{proof}

\begin{proof}[Proof of \cref{thm.promonad_theta}: identity]
The identity at $\lensob{c}$ has $\outp{f}=\id$ and $\inpt{f}=\eta(\varepsilon\otimes\id_{\inpt{c}})$.
\begin{align}
\phi_{\id}
&=\alpha\bigl(P(\id,\ev)\bigl(\lambda(\eta(\id_{H_c}),\eta(\varepsilon\otimes\id_{\inpt{c}}))\bigr)\bigr)
&&\text{definition}\\
&=\alpha\bigl(P(\id,\ev)\bigl(\eta(\id_{H_c}\otimes\varepsilon\otimes\id_{\inpt{c}})\bigr)\bigr)
&&\text{lax-unit (\cref{def.monoidal_promonad}(1))}\\
&=\alpha\bigl(\eta(\ev\circ(\varepsilon\otimes\id_{\inpt{c}}\otimes\id_{H_c}))\bigr)
&&\text{functoriality of }\eta\\
&=\ev\circ(\varepsilon\otimes\id_{\inpt{c}}\otimes\id_{H_c})
&&\alpha\text{-unit (\cref{def.P_alg}(1))}
\end{align}
Currying gives $\psi_{\id}=\varepsilon^*\colon\outp{c}\otimes H_c\to H_c$. Then $\Theta_m(\id)=(\id\otimes\varepsilon^*)\circ(\delta\otimes\id)=\id$ by the counit axiom.
\end{proof}

\begin{proof}[Proof of \cref{thm.promonad_theta}: composition]
Given $f\colon\lensob{c}\to\lensob{d}$ and $g\colon\lensob{d}\to\lensob{e}$, we must show $\Theta_m(g)\circ\Theta_m(f)=\Theta_m(g\circ f)$. The top-right boundary of the following diagram is $\Theta_m(g)\circ\Theta_m(f)$; the left-bottom boundary is $\Theta_m(g\circ f)$, with $\psi_{g\circ f}$ decomposed via \cref{lem.psi_compose}. We write juxtaposition for $\otimes$.
\[
\resizebox{\linewidth}{!}{%
\begin{tikzcd}[column sep=1.8em, row sep=1.8em, every label/.append style={font=\tiny}, font=\scriptsize, ampersand replacement=\&]
\outp c H_c
  \ar[r, "{\delta}"]
  \ar[d, "{\delta}"']
\& \outp c\outp c H_c \ar[r, "{\outp f}"] \ar[d, "{\delta}"']
\& \outp d\outp c H_c \ar[r, "{\psi_f}"] \ar[d, "{\delta}"']
\& \outp d H_d \ar[d, "{\delta}"]
\\
\outp c\outp c H_c \ar[d, "{\outp f}"'] \ar[r, "{\delta}"]
\& \outp c\outp c\outp c H_c \ar[r, "{\outp f\outp f}"] \ar[d, "{\outp f}"']
\& \outp d\outp d\outp c H_c \ar[r, "{\psi_f}"] \ar[d, "{\outp g}"']
\& \outp d\outp d H_d \ar[d, "{\outp g}"]
\\
\outp d\outp c H_c \ar[d, "{\outp g}"'] \ar[r, "{\delta}"]
\& \outp d\outp c\outp c H_c \ar[r, "{\outp g\outp f}"] \ar[d, "{\outp g}"']
\& \outp e\outp d\outp c H_c \ar[r, "{\psi_f}"] \ar[d, equal]
\& \outp e\outp d H_d \ar[d, equal]
\\
\outp e\outp c H_c \ar[r, "{\delta}"]
\& \outp e\outp c\outp c H_c \ar[r, "{\outp f}"]
\& \outp e\outp d\outp c H_c \ar[r, "{\psi_f}"]
\& \outp e\outp d H_d \ar[r, "{\psi_g}"]
\& \outp e H_e
\end{tikzcd}
}
\]
Every tile commutes by functoriality of $\otimes$ (operations on disjoint tensor factors), coassociativity of $\delta_{\outp{c}}$, or the comonoid-homomorphism property of $\outp{f}$.
\end{proof}

\bigskip
\hrule
\bigskip

\section{Reference: monad proof diagram}\label{page.ref_diagram}

The following diagram is reproduced from the main paper. It proves functoriality of $\Theta_z\colon\Lens{\cat{C}}^{\Fun{T}}\to\cat{C}$ for a strong monad $\Fun{T}$ and $\Fun{T}$-algebra $(z,\alpha)$. Write $H_X=\ihom{\inpt{X}}{z}$ and use juxtaposition for $\otimes$, $X$ for $\inpt{X}$.

\begin{landscape}
\thispagestyle{empty}
\begin{center}
\resizebox{\linewidth}{!}{%
\begin{tikzcd}[column sep=0.4em, row sep=1em, every label/.append style={font=\tiny}, font=\tiny, ampersand replacement=\&]
\outp c H_c
  \ar[r, "{\delta}"]
  \ar[d, "{\delta}"']
\& \outp c\outp c H_c \ar[r, "{\outp f}"] \ar[d, "{\delta}"']
\& \outp d\outp c H_c \ar[r, "{\coev_d}"] \ar[d, "{\delta}"']
\& \outp d[d,\outp c d H_c] \ar[r, "{\inpt f}"] \ar[d, "{\delta}"']
\& \outp d[d,\Fun T(c)H_c] \ar[r, "{\eta}"] \ar[d, "{\delta}"']
\& \outp d[d,\Fun T(c)\Fun T(H_c)] \ar[r, "{\tau}"] \ar[d, "{\delta}"']
\& \outp d[d,\Fun T(cH_c)] \ar[r, "{\Fun T(\ev)}"] \ar[d, "{\delta}"']
\& \outp d[d,\Fun T(z)] \ar[r, "{\alpha}"] \ar[d, "{\delta}"']
\& \outp d H_d \ar[d, "{\delta}"] \\
\outp c\outp c H_c \ar[d, "{\outp f}"'] \ar[r, "{\delta}"]
\& \outp c\outp c\outp c H_c \ar[r, "{\outp f\outp f}"] \ar[d, "{\outp f}"']
\& \outp d\outp d\outp c H_c \ar[r, "{\coev_d}"] \ar[d, "{\outp g}"']
\& \outp d\outp d[d,\outp c d H_c] \ar[r, "{\inpt f}"] \ar[d, "{\outp g}"']
\& \outp d\outp d[d,\Fun T(c)H_c] \ar[r, "{\eta}"] \ar[d, "{\outp g}"']
\& \outp d\outp d[d,\Fun T(c)\Fun T(H_c)] \ar[r, "{\tau}"] \ar[d, "{\outp g}"']
\& \outp d\outp d[d,\Fun T(cH_c)] \ar[r, "{\Fun T(\ev)}"] \ar[d, "{\outp g}"']
\& \outp d\outp d[d,\Fun T(z)] \ar[r, "{\alpha}"] \ar[d, "{\outp g}"']
\& \outp d\outp d H_d \ar[d, "{\outp g}"] \\
\outp d\outp c H_c \ar[d, "{\outp g}"'] \ar[r, "{\delta}"]
\& \outp d\outp c\outp c H_c \ar[r, "{\outp g\outp f}"]
\& \outp e\outp d\outp c H_c \ar[r, "{\coev_d}"] \ar[d, "{\coev_e}"']
\& \outp e\outp d[d,\outp c d H_c] \ar[r, "{\inpt f}"] \ar[d, "{\coev_e}"']
\& \outp e\outp d[d,\Fun T(c)H_c] \ar[r, "{\eta}"] \ar[d, "{\coev_e}"']
\& \outp e\outp d[d,\Fun T(c)\Fun T(H_c)] \ar[r, "{\tau}"] \ar[d, "{\coev_e}"']
\& \outp e\outp d[d,\Fun T(cH_c)] \ar[r, "{\Fun T(\ev)}"] \ar[d, "{\coev_e}"']
\& \outp e\outp d[d,\Fun T(z)] \ar[r, "{\alpha}"] \ar[d, "{\coev_e}"']
\& \outp e\outp d H_d \ar[d, "{\coev_e}"] \\
\outp e\outp c H_c \ar[dd, "{\coev_e}"']
\&\& \outp e[e,\outp d e\outp c H_c] \ar[r, "{\coev_d}"] \ar[d, "{\inpt g}"']
\& \outp e[e,\outp d e[d,\outp c d H_c]] \ar[r, "{\inpt f}"] \ar[d, "{\inpt g}"']
\& \outp e[e,\outp d e[d,\Fun T(c)H_c]] \ar[r, "{\eta}"] \ar[d, "{\inpt g}"']
\& \outp e[e,\outp d e[d,\Fun T(c)\Fun T(H_c)]] \ar[r, "{\tau}"] \ar[d, "{\inpt g}"']
\& \outp e[e,\outp d e[d,\Fun T(cH_c)]] \ar[r, "{\Fun T(\ev)}"] \ar[d, "{\inpt g}"']
\& \outp e[e,\outp d e[d,\Fun T(z)]] \ar[r, "{\alpha}"] \ar[d, "{\inpt g}"']
\& \outp e[e,\outp d e H_d] \ar[d, "{\inpt g}"] \\
\& \outp e[e,\Fun T(\outp c)\Fun T(d)H_c] \ar[d, "{\eta}"'] \ar[dddd, out=200, in=150, equal]
\& \outp e[e,\outp c\Fun T(d)H_c] \ar[r, "{\coev_d}"] \ar[d, "{\eta}"] \ar[l, "{\eta}"]
\& \outp e[e,\Fun T(d)[d,\outp c d H_c]] \ar[r, "{\inpt f}"] \ar[d, "{\eta}"']
\& \outp e[e,\Fun T(d)[d,\Fun T(c)H_c]] \ar[r, "{\eta}"] \ar[d, "{\eta}"']
\& \outp e[e,\Fun T(d)[d,\Fun T(c)\Fun T(H_c)]] \ar[r, "{\tau}"] \ar[d, "{\eta}"']
\& \outp e[e,\Fun T(d)[d,\Fun T(cH_c)]] \ar[r, "{\Fun T(\ev)}"] \ar[d, "{\eta}"']
\& \outp e[e,\Fun T(d)[d,\Fun T(z)]] \ar[r, "{\alpha}"] \ar[d, "{\eta}"']
\& \outp e[e,\Fun T(d)H_d] \ar[d, "{\eta}"] \\
\outp e[e,\outp c e H_c] \ar[d, "{\delta}"']
\& \outp e[e,\Fun T(\outp c)\Fun T(d)\Fun T(H_c)] \ar[r, "{\tau}"]
\& \outp e[e,\Fun T(\outp c d)\Fun T(H_c)] \ar[r, "{\Fun T(\coev_d)}"] \ar[d, "{\tau}"']
\& \outp e[e,\Fun T(d)\Fun T([d,\outp c d H_c])] \ar[r, "{\Fun T(\inpt f)}"] \ar[d, "{\tau}"']
\& \outp e[e,\Fun T(d)\Fun T([d,\Fun T(c)H_c])] \ar[r, "{\Fun T(\eta)}"] \ar[d, "{\tau}"']
\& \outp e[e,\Fun T(d)\Fun T([d,\Fun T(c)\Fun T(H_c)])] \ar[r, "{\Fun T(\tau)}"] \ar[d, "{\tau}"']
\& \outp e[e,\Fun T(d)\Fun T([d,\Fun T(cH_c)])] \ar[r, "{\Fun T(\ev)}"] \ar[d, "{\tau}"']
\& \outp e[e,\Fun T(d)\Fun T([d,\Fun T(z)])] \ar[r, "{\Fun T(\alpha)}"] \ar[d, "{\tau}"']
\& \outp e[e,\Fun T(d)\Fun T(H_d)] \ar[d, "{\tau}"] \\
\outp e[e,\outp c\outp c e H_c] \ar[d, "{\outp f}"']
\&\& \outp e[e,\Fun T(\outp c d H_c)] \ar[r, "{\Fun T(\coev_d)}"] \ar[dr, equal]
\& \outp e[e,\Fun T(d[d,\outp c d H_c])] \ar[r, "{\Fun T(\inpt f)}"] \ar[d, "{\Fun T(\ev)}"']
\& \outp e[e,\Fun T(d[d,\Fun T(c)H_c])] \ar[r, "{\Fun T(\eta)}"] \ar[d, "{\Fun T(\ev)}"']
\& \outp e[e,\Fun T(d[d,\Fun T(c)\Fun T(H_c)])] \ar[r, "{\Fun T(\tau)}"] \ar[d, "{\Fun T(\ev)}"']
\& \outp e[e,\Fun T(d[d,\Fun T(cH_c)])] \ar[r, "{\Fun T(\ev)}"] \ar[d, "{\Fun T(\ev)}"']
\& \outp e[e,\Fun T(d[d,\Fun T(z)])] \ar[r, "{\Fun T(\alpha)}"] \ar[d, "{\Fun T(\ev)}"']
\& \outp e[e,\Fun T(dH_d)] \ar[d, "{\Fun T(\ev)}"] \\
\outp e[e,\outp d e\outp c H_c] \ar[d, "{\inpt g}"']
\&\&
\& \outp e[e,\Fun T(\outp c d H_c)] \ar[r, "{\Fun T(\inpt f)}"]
\& \outp e[e,\Fun T(\Fun T(c)H_c)] \ar[r, "{\Fun T(\eta)}"]
\& \outp e[e,\Fun T(\Fun T(c)\Fun T(H_c))] \ar[r, "{\Fun T(\tau)}"]
\& \outp e[e,\Fun T(\Fun T(cH_c))] \ar[r, "{\Fun T(\ev)}"] \ar[d, "{\mu}"']
\& \outp e[e,\Fun T(\Fun T(z))] \ar[r, "{\Fun T(\alpha)}"] \ar[d, "{\mu}"']
\& \outp e[e,\Fun T(z)] \ar[d, "{\alpha}"] \\
\outp e[e,\outp c\Fun T(d)H_c] \ar[r, "{\eta}"]
\& \outp e[e,\Fun T(\outp c)\Fun T(d)H_c] \ar[r, "{\tau}"]
\& \outp e[e,\Fun T(\outp c d)H_c] \ar[r, "{\Fun T(\inpt f)}"]
\& \outp e[e,\Fun T(\Fun T(c))H_c] \ar[r, "{\mu}"]
\& \outp e[e,\Fun T(c)H_c] \ar[r, "{\eta}"]
\& \outp e[e,\Fun T(c)\Fun T(H_c)] \ar[r, "{\tau}"]
\& \outp e[e,\Fun T(cH_c)] \ar[r, "{\Fun T(\ev)}"]
\& \outp e[e,\Fun T(z)] \ar[r, "{\alpha}"]
\& \outp e H_e
\end{tikzcd}
}
\end{center}

\medskip

The curved equality on the left splits the open region into two parts, and it suffices to show each commutes. The left-hand square below shows that the two maps $\outp d\otimes\outp c\otimes[\inpt c,z]\to\outp e\otimes[\inpt e,\Fun T(\outp c)\otimes \Fun T(d)\otimes[\inpt c,z]]$ agree, and the right-hand square shows that the two maps $\outp e\otimes[\inpt e,\Fun T(\outp c)\otimes\Fun T(d)\otimes [\inpt c,z]]\to\outp e\otimes [\inpt e,\Fun T(\inpt c\otimes [\inpt c,z])]$ agree.

\noindent\begin{minipage}{0.48\linewidth}
\[
\resizebox{\linewidth}{!}{%
\begin{tikzcd}[column sep=1em, row sep=1em, every label/.append style={font=\tiny}, font=\scriptsize, ampersand replacement=\&]
\outp d\outp c H_c
  \ar[r, "{\delta}"]
  \ar[d, "{\outp g}"']
\& \outp d\outp c\outp c H_c \ar[r, "{\outp g}"] \ar[d, "{\outp g}"']
\& \outp e\outp c\outp c H_c \ar[r, "{\outp f}"] \ar[d, "{\outp f}"']
\& \outp e\outp d\outp c H_c \ar[d, "{\coev_e}"] \\
\outp e\outp c H_c \ar[r, "{\delta}"] \ar[d, "{\coev_e}"']
\& \outp e\outp c\outp c H_c \ar[r, "{\outp f}"] \ar[d, "{\coev_e}"']
\& \outp e\outp d\outp c H_c \ar[r, "{\coev_e}"] \ar[d, "{\coev_e}"']
\& \outp e[e,\outp d e\outp c H_c] \ar[d, "{\inpt g}"] \\
\outp e[e,\outp c e H_c] \ar[r, "{\delta}"] \ar[d, "{\delta}"']
\& \outp e[e,\outp c\outp c e H_c] \ar[r, "{\outp f}"] \ar[d, "{\outp f}"']
\& \outp e[e,\outp d e\outp c H_c] \ar[r, "{\inpt g}"] \ar[d, "{\inpt g}"']
\& \outp e[e,\outp c\Fun T(d)H_c] \ar[d, "{\eta}"] \\
\outp e[e,\outp c\outp c e H_c] \ar[r, "{\outp f}"]
\& \outp e[e,\outp d e\outp c H_c] \ar[r, "{\inpt g}"]
\& \outp e[e,\outp c\Fun T(d)H_c] \ar[r, "{\eta}"]
\& \outp e[e,\Fun T(\outp c)\Fun T(d)H_c]
\end{tikzcd}
}
\]
\end{minipage}\hfill
\begin{minipage}{0.5\linewidth}
\[
\resizebox{\linewidth}{!}{%
\begin{tikzcd}[column sep=1em, row sep=1em, every label/.append style={font=\tiny}, font=\scriptsize, ampersand replacement=\&]
\outp e[e,\Fun T(\outp c)\Fun T(d)H_c]
  \ar[r, "{\eta}"]
  \ar[d, "{\tau}"']
\& \outp e[e,\Fun T(\outp c)\Fun T(d)\Fun T(H_c)] \ar[r, "{\tau}"] \ar[d, "{\tau}"']
\& \outp e[e,\Fun T(\outp c d)\Fun T(H_c)] \ar[r, "{\tau}"]
\& \outp e[e,\Fun T(\outp c d H_c)] \ar[d, "{\Fun T(\inpt f)}"] \\
\outp e[e,\Fun T(\outp c d)H_c] \ar[r, "{\eta}"] \ar[dd, "{\Fun T(\inpt f)}"']
\& \outp e[e,\Fun T(\outp c d)\Fun T(H_c)] \ar[urr, "{\tau}"] \ar[d, "{\Fun T(\inpt f)}"']
\&\& \outp e[e,\Fun T(\Fun T(c)H_c)] \ar[d, "{\Fun T(\eta)}"] \\
\& \outp e[e,\Fun T(\Fun T(c))\Fun T(H_c)] \ar[urr, "{\tau}"] \ar[dd, "{\mu}"'] \ar[r, "{\Fun T(\eta)}"']
\& \outp e[e,\Fun T(\Fun T(c))\Fun T(\Fun T(H_c))] \ar[r, "{\tau}"] \ar[d, "{\mu}"']
\& \outp e[e,\Fun T(\Fun T(c)\Fun T(H_c))] \ar[d, "{\Fun T(\tau)}"] \\
\outp e[e,\Fun T(\Fun T(c))H_c] \ar[ur, "{\eta}"] \ar[d, "{\mu}"']
\&\& \outp e[e,\Fun T(c)\Fun T(\Fun T(H_c))] \ar[dl, "{\mu}"]
\& \outp e[e,\Fun T(\Fun T(cH_c))] \ar[d, "{\mu}"] \\
\outp e[e,\Fun T(c)H_c] \ar[r, "{\eta}"]
\& \outp e[e,\Fun T(c)\Fun T(H_c)] \ar[rr, "{\tau}"]
\&\& \outp e[e,\Fun T(cH_c)]
\end{tikzcd}
}
\]
\end{minipage}
\end{landscape}

\noindent To obtain the promonad proof, collapse columns 4--8 into a single column $\ihom{d}{\phi_f}$. The top three rows (comonoid axioms) remain unchanged. Rows 4--9 collapse via the four steps of \cref{sec.comp}: the algebra axiom replaces the $\mu$/$\Fun{T}\alpha$ square (bottom-right of the diagram), Tambara naturality replaces the $\sigma$-naturality squares (rows 5--6), and naturality of $\alpha$ replaces the naturality-of-$a$ squares (right column).

\end{document}
